\documentclass[a4paper,11pt]{jsarticle}
\usepackage{amsmath,amssymb,amsthm,mathtools}
\usepackage{geometry}
\geometry{margin=25mm}
\newtheorem{definition}{定義}[section]
\newtheorem{remark}[definition]{注意}
\newcommand{\Solver}{\mathcal{S}}
\newcommand{\Ax}{S_{\mathrm{ax}}}
\newcommand{\Th}{\mathrm{Th}}
\newcommand{\Zset}{\mathcal{Z}}
\newcommand{\Mod}{\mathrm{Mod}}

\title{ECA_8\\$\mathcal{Z}$ におけるモデル分岐と連続体仮説}
\author{}
\date{}

\begin{document}
\maketitle

\begin{abstract}
本稿では、ECA_7 において再定義した
\[
\mathcal{Z}=
\left\{s\in\Solver_I\mid
\mathcal{T}(s)\cong_{\mathrm{Th}}\mathrm{ZFC}\right\}
\]
を前提として、ECA_5 で記述されていた連続体仮説（CH）の
「Solver に依存する」という記述を、理論レベルとモデルレベルに
分離して再構成する。

特に、異なる Solver が ZFC と理論同型な理論を表現することと、
それらの理論について CH または $\neg\mathrm{CH}$ を満たすモデルを
考えることを区別する。

本稿の目的は CH の独立性そのものを新たに証明することではない。
また、$\mathcal{Z}$ の濃度や連続体濃度 $2^{\aleph_0}$ を
本稿のモデル分岐から導出するものでもない。
\end{abstract}

\section{本稿の位置付け}

ECA_7 では、ECA に適合する Solver 全体を
\[
\Solver_I=
\left\{s=(\mathcal{A}_s,\mathcal{R}_s)\mid
\mathcal{A}_s\in I\right\}
\]
とし、理論写像を
\[
\mathcal{T}:\Solver_I\longrightarrow\Th
\]
とした。

さらに、
\[
\mathcal{Z}=
\left\{s\in\Solver_I\mid
\mathcal{T}(s)\cong_{\mathrm{Th}}\mathrm{ZFC}\right\}
\]
を定義した。

したがって、$s\in\mathcal{Z}$ は「Solver $s$ 自体が ZFC である」
という意味ではなく、Solver $s$ が表現する理論
$\mathcal{T}(s)$ が ZFC と理論同型であることを意味する。

\section{ECA_5 における従来の CH 記述}

ECA_5 では、
\[
\exists S_1,S_2\in\mathcal{Z}\quad
\left(
\mathcal{T}(S_1)\models\mathrm{CH}
\land
\mathcal{T}(S_2)\models\neg\mathrm{CH}
\right)
\]
という形で CH の分岐を記述していた。

しかし、ECA_7 で理論とモデルを明確に区別したため、
この式は理論レベルの主張として読まれる可能性がある。
本稿では、ECA_5 の意図していた分岐をモデルを明示した形へ改める。

\section{モデルの導入}

\begin{definition}[モデル]
各理論 $T\in\Th$ に対して、そのモデル全体を
\[
\Mod(T)
\]
と表す。

\[
M\in\Mod(T)
\]
とは、
\[
M\models T
\]
を満たす構造 $M$ であることを表す。
\end{definition}

したがって Solver $s$ に対しては、
\[
M\in\Mod(\mathcal{T}(s))
\]
によって $\mathcal{T}(s)$ のモデルを表す。

\subsection{理論とモデルの論理関係}

\[
\mathcal{T}(s)\models\mathrm{CH}
\]
と
\[
\exists M\in\Mod(\mathcal{T}(s))\quad
M\models\mathrm{CH}
\]
は異なる主張である。

前者は CH が理論 $\mathcal{T}(s)$ の論理的帰結であることを表す。
後者は $\mathcal{T}(s)$ のモデルのうち CH を満たすものが存在する
ことを表す。

同様に、
\[
\exists M\in\Mod(\mathcal{T}(s))\quad
M\models\neg\mathrm{CH}
\]
は $\neg\mathrm{CH}$ を満たすモデルの存在を表す。

\section{CH と $\neg$CH のモデル分岐}

\begin{definition}[CH モデル]
$s\in\mathcal{Z}$ に対して、
\[
M_{\mathrm{CH}}\in\Mod(\mathcal{T}(s)),
\qquad
M_{\mathrm{CH}}\models\mathrm{CH}
\]
を満たすモデルを CH モデルと呼ぶ。
\end{definition}

\begin{definition}[$\neg$CH モデル]
$s\in\mathcal{Z}$ に対して、
\[
M_{\neg\mathrm{CH}}\in\Mod(\mathcal{T}(s)),
\qquad
M_{\neg\mathrm{CH}}\models\neg\mathrm{CH}
\]
を満たすモデルを $\neg$CH モデルと呼ぶ。
\end{definition}

異なる Solver を明示したモデル分岐は、次のように記述できる。
\[
\boxed{
\begin{aligned}
\exists s_1,s_2\in\mathcal{Z},\;
\exists M_1,M_2\quad
&
\bigl(
M_1\in\Mod(\mathcal{T}(s_1))
\land M_1\models\mathrm{CH}
\\
&{}\land
M_2\in\Mod(\mathcal{T}(s_2))
\land M_2\models\neg\mathrm{CH}
\bigr).
\end{aligned}}
\]

ただし、これは「異なる Solver を用いたモデル分岐」を表す条件であり、
ECA の定義から自動的に導かれるものではない。

\section{同一 Solver におけるモデル分岐}

モデル分岐と Solver の分岐は同一ではない。

一つの Solver $s\in\mathcal{Z}$ に対して
\[
M_1,M_2\in\Mod(\mathcal{T}(s))
\]
かつ
\[
M_1\models\mathrm{CH},
\qquad
M_2\models\neg\mathrm{CH}
\]
という状況を考えることも形式上は可能である。

この場合、
\[
\boxed{\text{Solver の分岐}\neq\text{モデルの分岐}}
\]
となる。

すなわち、モデルの違いを表すために
\[
s_1\neq s_2
\]
を必ずしも要求しない。

\section{異なる Solver を用いた分岐}

一方、ECA の Solver 構造そのものの違いを併せて記述する場合には、
\[
s_1,s_2\in\mathcal{Z},
\qquad s_1\neq s_2
\]
を考え、
\[
M_1\in\Mod(\mathcal{T}(s_1)),
\qquad
M_2\in\Mod(\mathcal{T}(s_2))
\]
とする。

さらに、
\[
M_1\models\mathrm{CH},
\qquad
M_2\models\neg\mathrm{CH}
\]
を追加条件として置く。

ここで重要なのは、
\[
s_1\neq s_2
\]
だけから CH と $\neg$CH の違いが導かれるわけではないことである。

また、
\[
\mathcal{T}(s_1)\cong_{\mathrm{Th}}\mathcal{T}(s_2)
\]
だけから、二つのモデルにおける CH の真偽が異なるとも限らない。

したがって、モデル分岐は明示的な追加条件として扱う。

\section{ZFC との理論同型と CH}

$\mathcal{Z}$ の定義から、
\[
s\in\mathcal{Z}
\Longrightarrow
\mathcal{T}(s)\cong_{\mathrm{Th}}\mathrm{ZFC}
\]
である。

しかし、この条件から
\[
\mathcal{T}(s)\models\mathrm{CH}
\]
または
\[
\mathcal{T}(s)\models\neg\mathrm{CH}
\]
のいずれかが直ちに導かれるわけではない。

理論同型は理論としての構造を対応させる条件であり、
特定のモデルを一つ選択する条件ではない。

したがって、
\[
\boxed{
\mathcal{T}(s)\cong_{\mathrm{Th}}\mathrm{ZFC}
}
\]
と
\[
\boxed{
M\models\mathrm{CH}
}
\]
は異なる層の条件である。

\section{ECA_5 の CH 記述の再構成}

ECA_5 の記述を ECA_7 の定義に合わせるなら、
「CH の真偽は Solver 集合の選択に依存する」と直接断定するのではなく、
まず
\[
s\in\mathcal{Z}
\]
を ZFC と理論同型な理論を表現する Solver とし、
そのモデル
\[
M\in\Mod(\mathcal{T}(s))
\]
を考える。

その上で、
\[
M\models\mathrm{CH}
\]
または
\[
M\models\neg\mathrm{CH}
\]
というモデルレベルの分岐を記述する。

異なる Solver を比較する場合には、
\[
s_1,s_2\in\mathcal{Z}
\]
に対して
\[
M_1\in\Mod(\mathcal{T}(s_1)),
\qquad
M_2\in\Mod(\mathcal{T}(s_2))
\]
を取り、
\[
M_1\models\mathrm{CH},
\qquad
M_2\models\neg\mathrm{CH}
\]
という関係を明示する。

これにより、
\[
\mathcal{Z}
\]
の定義と、その元に対するモデルの分岐を混同しない。

\section{モデル分岐と $\mathcal{Z}$ の濃度}

本稿で導入したモデル分岐から、
\[
|\mathcal{Z}|
\]
の値を決定することはできない。

特に、
\[
\exists s_1,s_2\in\mathcal{Z}
\]
という存在条件が成立したとしても、
それだけから
\[
|\mathcal{Z}|\geq2^{\aleph_0}
\]
は導かれない。

同様に、二つのモデルについて
\[
M_1\models\mathrm{CH},
\qquad
M_2\models\neg\mathrm{CH}
\]
が成立しても、
\[
|\mathcal{Z}|=2^{\aleph_0}
\]
は導かれない。

したがって、$\mathcal{Z}$ の濃度については、
モデル分岐とは独立に、Solver の構成と写像を用いて別途検討する。

\section{本稿で確定した関係}

以上により、ECA における関係を
\[
\boxed{
s
\longrightarrow
\mathcal{T}(s)
\longrightarrow
M
}
\]
と整理する。

ここで、
\[
s\in\Solver_I,\qquad
\mathcal{T}(s)\in\Th,\qquad
M\in\Mod(\mathcal{T}(s))
\]
である。

また、
\[
s\in\mathcal{Z}
\iff
\mathcal{T}(s)\cong_{\mathrm{Th}}\mathrm{ZFC}
\]
は理論レベルの条件であり、
\[
M\models\mathrm{CH}
\quad\text{または}\quad
M\models\neg\mathrm{CH}
\]
はモデルレベルの条件である。

したがって、
\[
\boxed{
\begin{array}{c}
\text{Solver}\\
\downarrow\\
\text{理論}\\
\downarrow\\
\text{モデル}\\
\downarrow\\
\mathrm{CH}/\neg\mathrm{CH}
\end{array}}
\]
という階層を維持する。

\section{後続する議論}

本稿を基礎として、後続文書では次を個別に検討できる。

\begin{enumerate}
\item $\mathcal{Z}$ に含まれる Solver の具体的構成。
\item 異なる Solver が $\mathcal{Z}$ の異なる元となる条件。
\item $\mathfrak{A}\restriction_{\mathcal{Z}}$ と
$\mathcal{T}\restriction_{\mathcal{Z}}$ の関係。
\item $\mathcal{Z}$ の濃度と $2^{\aleph_0}$ との関係。
\item 必要に応じて CH/$\neg$CH を満たすモデルの構成条件。
\end{enumerate}

これらは本稿の定義から自動的に導かれるものではなく、
それぞれ追加の構成または証明を必要とする。

\end{document}
